% GENERATED, edit the French unit or the translation table
\mxnewpage\section*{Solutions EV1-T1}\label{EV1-T1-CORR}
\textbf{MCQ:} 1 B, 2 C, 3 A, 4 C, 5 D, 6 D, 7 B, 8 D, 9 D, 10 B, 11 B, 12 A. Each correct answer: 0.5 point; total 6.
\paragraph{Checks of the fluency questions.} Order of operations; sum of fractions; subtraction of exponents; multiplying coefficient 1.2; division by 0.8; product 1.25 times 0.8; solving an affine equation; vertical increment divided by horizontal; mean; angle divided by 360; complement to 1; substitution in brackets.
\subsection*{Exercise 1 -- A book collection}
\begin{enumerate}
\item $u_n=120+18n$, arithmetic with common difference 18. \textit{Marking: 1 point. Result and basic justification expected; half a point possible for a correct but unfinished approach.}
\item $u_8=264$; nine terms, sum $\dfrac{9(120+264)}{2}=1728$. This sum accumulates monthly states, not distinct books. \textit{Marking: 2 points. Method or justification: 1; calculation and conclusion: 1.}
\item $120+18n>300\iff n>10$. First index 11; $u_{10}=300$, $u_{11}=318$. \textit{Marking: 2 points. Method or justification: 1; calculation and conclusion: 1.}
\end{enumerate}
\subsection*{Exercise 2 -- Area and choice of form}
\begin{enumerate}
\item $A=-x^2+14x=49-(x-7)^2$. \textit{Marking: 1 point. Result and basic justification expected; half a point possible for a correct but unfinished approach.}
\item $A\leq49$, equality at 7. \textit{Marking: 2 points. Method or justification: 1; calculation and conclusion: 1.}
\item $49-(x-7)^2\geq40\iff(x-7)^2\leq9$, so $[4;10]$. \textit{Marking: 2 points. Method or justification: 1; calculation and conclusion: 1.}
\end{enumerate}
\subsection*{Exercise 3 -- A right angle}
\begin{enumerate}
\item $\vect{AB}=(3;1)$, $\vect{AC}=(-1;3)$. \textit{Marking: 1 point. Result and basic justification expected; half a point possible for a correct but unfinished approach.}
\item Product $-3+3=0$. \textit{Marking: 1 point. Result and basic justification expected; half a point possible for a correct but unfinished approach.}
\item $BC=\sqrt{(-4)^2+2^2}=2\sqrt5$; perpendicular sides of norm $\sqrt{10}$, area 5. \textit{Marking: 2 points. Method or justification: 1; calculation and conclusion: 1.}
\end{enumerate}
\paragraph{Calibration.} Target duration: 20--25 min for the fluency questions, 80--85 min for the exercises, 10--15 min for proofreading. These durations are editorial estimates to be measured in class. Accept any correct method; do not penalise twice a numerical error that is propagated consistently. The proposed marking scheme is not a decision of the team to harmonise marking.
